\section{Integrated rigidity in dimensions two and three}
\label{uni-sec:lowdim}

The linearized equation is not the only way to use equality of
determinants.  In low dimension a quadratic difference has a definite sign,
and its integral is fixed by exactness or by a sign on the torsion.

\begin{theorem}\label{uni-thm:sobolev}
Let $n\in\{2,3\}$ and let $u,v\in W^{2,n}(X)$.  Assume
$A,B\geq0$ almost everywhere and
\[
  A^n=B^n=f\,\mathrm{d}V,\qquad f>0\quad\hbox{almost everywhere}.
\]
For $n=2$ impose no differential condition on $\omega$.  For
$n=3$, assume
\begin{equation}
  \ddc\omega\wedge\mathrm{i}\eta\wedge\bar\eta\leq0
  \label{uni-eq:threefold-sign}
\end{equation}
for every $(1,0)$-covector $\eta$.  Then
$\ddc(u-v)=0$ distributionally and $u-v$ is almost everywhere constant.
For upper semicontinuous $\omega$-plurisubharmonic representatives the
equality holds pointwise.
\end{theorem}

\begin{proof}
Put
\[
 \begin{gathered}
 w=u-v,\qquad s=(u+v)/2,\qquad D=\ddc w=A-B,\\
 C=(A+B)/2=\omega+\ddc s.
 \end{gathered}
\]
The coefficients of $C,D$ lie in $L^n$; their top products are
integrable by H\"older's inequality.  At almost every point the common
determinant is positive, so $A,B$ are positive definite.  Diagonalize
them simultaneously by congruence, writing
$A=\sum e_j$, $B=\sum\lambda_j e_j$, where
$e_j=\mathrm{i}\theta_j\wedge\bar\theta_j$, $\lambda_j>0$, and
$\prod\lambda_j=1$.

In dimension two,
\begin{equation}\label{uni-eq:surface-sign}
 D^2=A^2+B^2-2A\wedge B
    =2(2-\lambda_1-\lambda_2)e_1\wedge e_2\leq0,
\end{equation}
with equality precisely when $A=B$.  In dimension three,
\begin{equation}\label{uni-eq:threefold-algebra}
 2C\wedge D^2=A^3+B^3-A^2\wedge B-A\wedge B^2\leq0.
\end{equation}
Indeed, with $E=e_1\wedge e_2\wedge e_3$, the four terms on the right
are
\[
 6E,\qquad 6E,\qquad 2\Bigl(\sum_j\lambda_j\Bigr)E,
 \qquad 2\Bigl(\sum_{i<j}\lambda_i\lambda_j\Bigr)E.
\]
Each sum is at least three, and equality again forces all
$\lambda_j=1$.

Choose smooth real $w_k\to w$ in $W^{2,n}$ and set
$D_k=\ddc w_k$.  Such approximations follow by a partition of unity,
local extension of compactly supported pieces, and convolution.  They
are not required to preserve semipositivity.  For $n=2$, Stokes gives
$\int_XD_k^2=0$, since
$D_k^2=\mathrm{d}(\dc w_k\wedge D_k)$.  Moreover
\[
 \|D_k^2-D^2\|_{L^1}
 \leq C\|D_k-D\|_{L^2}(\|D_k\|_{L^2}+\|D\|_{L^2})\longrightarrow0.
\]
Thus $\int_XD^2=0$, and the nonpositive integrand in
\eqref{uni-eq:surface-sign} vanishes almost everywhere.  Its strict equality
case gives $D=0$ almost everywhere.

For $n=3$, also choose smooth $s_k\to s$ in $W^{2,3}$, and put
$C_k=\omega+\ddc s_k$.  The needed identity is
\begin{equation}\label{uni-eq:torsion-integral}
 \int_X C_k\wedge D_k^2
 =\int_X\omega\wedge D_k^2
 =-\int_X\mathrm{d}w_k\wedge \dc w_k\wedge\ddc\omega.
\end{equation}
Here the first equality follows because
$\ddc s_k\wedge D_k^2=\mathrm{d}(\dc s_k\wedge D_k^2)$.
To check the sign in the second equality, put
$J_k=\mathrm{d}w_k\wedge \dc w_k$.  Using
$\dc \mathrm{d}=-\ddc $ and $(\dc )^2=0$,
\[
 \dc J_k=-D_k\wedge \dc w_k,\qquad \ddc J_k=-D_k^2.
\]
Stokes, first for $\mathrm{d}(\omega\wedge \dc J_k)$ and then for
$\dc (\mathrm{d}\omega\wedge J_k)$, gives
\begin{align*}
 \int_X\omega\wedge\ddc J_k&=-\int_X\mathrm{d}\omega\wedge \dc J_k,\\
 0&=-\int_X\ddc\omega\wedge J_k-\int_X\mathrm{d}\omega\wedge \dc J_k.
\end{align*}
Combining them yields \eqref{uni-eq:torsion-integral}.  Thus no unrecorded
first-order torsion term is discarded.

Strong $L^3$ convergence implies convergence of all three-factor
products in $L^1$.  For example,
\[
 \|(C_k-C)\wedge D_k^2\|_1
 \leq C\|C_k-C\|_3\|D_k\|_3^2\longrightarrow0,
\]
and
\[
 \|C\wedge(D_k^2-D^2)\|_1
 \leq C\|C\|_3\|D_k-D\|_3(\|D_k\|_3+\|D\|_3)\longrightarrow0.
\]
Also $\mathrm{d}w_k\wedge \dc w_k\to \mathrm{d}w\wedge \dc w$ in $L^{3/2}$, so the
right side of \eqref{uni-eq:torsion-integral} converges.  Consequently
\begin{equation}\label{uni-eq:torsion-limit}
 \int_XC\wedge D^2=-\int_X\mathrm{d}w\wedge \dc w\wedge\ddc\omega\geq0.
\end{equation}
The last sign uses
$\mathrm{d}w\wedge \dc w=(\mathrm{i}/(2\pi))\partial w\wedge\bar\partial w$
and \eqref{uni-eq:threefold-sign}.  The left integrand is nonpositive by
\eqref{uni-eq:threefold-algebra}, so it vanishes almost everywhere.  The
positive-density equality case gives $D=0$ almost everywhere.

In both dimensions $\ddc w=0$ as a distribution.  Taking the complex
trace in a chart gives the distributional Euclidean equation
$\Delta w=0$.  To recall why a smooth representative exists, convolve
locally with smooth radial kernels.  The convolutions are harmonic on
smaller balls and converge to $w$ in $L^1$.  Their mean-value property
against a fixed smaller radial kernel passes to the limit, identifying
$w$ almost everywhere with a smooth harmonic convolution there.  These
representatives agree on overlaps and give a global smooth harmonic
representative.  Its attained maximum on compact connected $X$ and the
local mean-value property make it constant.

If the original potentials are upper semicontinuous
$\omega$-plurisubharmonic representatives, each equals the limit of its
local ball averages: this follows from the submean inequality and upper
semicontinuity after adding a smooth local function.  The almost
everywhere constant-difference equality therefore holds pointwise.  Equal
mean normalizations give equality almost everywhere in the Sobolev class;
equal supremum normalizations give equality for the plurisubharmonic
representatives.
\end{proof}

\begin{corollary}\label{uni-cor:dense-lowdim}
For $C^2$ potentials in dimensions two and three, the corresponding
conclusion remains valid if the common continuous density is positive
merely on a dense set.  In dimension three retain
\eqref{uni-eq:threefold-sign}; in particular $\ddc\omega=0$ is sufficient.
\end{corollary}

\begin{proof}
At a point where the common top power vanishes, the algebraic expressions
are still nonpositive: in dimension two $D^2=-2A\wedge B$, and in
dimension three
$2C\wedge D^2=-A^2\wedge B-A\wedge B^2$.
The same integral identities force these continuous expressions to vanish
everywhere.  Their strict equality cases give $D=0$ on $\{f>0\}$.
Density and continuity then give $D=0$ globally.  No equality-case
assertion at a degenerate point is needed.
\end{proof}


The Hopf threefold of \cref{chap:hopf} illustrates the distinction between this uniqueness argument and a mass bound. Its metric satisfies $\ddc\omega=-\alpha^2$ with $\alpha\geq0$, so it meets the sign condition in \eqref{uni-eq:threefold-sign}. The Sobolev uniqueness theorem therefore applies on this example, even though its upper Monge--Amp\`ere volume is infinite.
